\subsection{问题一的描述与分析}
本问在空气--外延层--衬底三层结构中，仅考虑表面反射光与第一束内部反射光，建立外延层厚度与红外反射光谱的定量关系。核心是确定两束光的光程差与传播相位，引入SiC折射率的色散与载流子响应，并由Fresnel界面系数完成相干叠加，得到以厚度和材料光学参数为输入、非偏振反射率为输出的双光束理论模型。

图\ref{fig:q1-flow}给出问题一从几何光路、材料色散到反射率模型的总体建模流程。

\begin{figure}[!htbp]
\centering
\includegraphics[width=\textwidth]{\ProjectRoot/modules/20_q1/figures/editable/origin_templates/reference_png/q1_flow_reference.png}
\caption{问题一Sellmeier--Drude--Fresnel双光束模型流程。}
\label{fig:q1-flow}
\end{figure}

\subsection{Sellmeier--Drude--Fresnel双光束模型}
依据薄膜干涉与分层介质传播理论\cite{bornwolf,yeh}，本文将几何传播相位、SiC材料色散和Fresnel界面响应统一到复振幅叠加框架，建立Sellmeier--Drude--Fresnel双光束测厚模型。多入射角Fresnel响应常用于联合表征薄膜厚度与光学常数\cite{politano2023}，但本文仅以此作为方法背景，具体材料关系采用下述SiC实验与计算文献。

\subsubsection{光程差与传播相位}
图\ref{fig:q1-schematic}为空气--外延层--衬底结构中的双光束干涉光路示意图。设空气、外延层和衬底折射率分别为$n_1,n_2,n_3$，外延层厚度为$d$，入射角为$i$，外延层折射角为$\gamma$。数值计算中取空气折射率$n_1=1.0003$。

\begin{figure}[!htbp]
\centering
\includegraphics[width=.70\textwidth]{\ProjectRoot/modules/20_q1/figures/q1_interference_schematic.png}
\caption{空气--外延层--衬底结构中的双光束干涉光路。}
\label{fig:q1-schematic}
\end{figure}

由分层介质中的Snell折射关系\cite{bornwolf,yeh}有
\begin{equation}
n_1\sin i=n_2\sin\gamma.
\label{eq:q1-snell}
\end{equation}
内部反射光在外延层中下行和上行的几何路程均为$d/\cos\gamma$。将两束出射光比较到同一等相面，表面反射光在空气中的等效补偿路程为$2d\tan\gamma\sin i$，因此
\[
\Delta
=2n_2\frac{d}{\cos\gamma}-2n_1d\tan\gamma\sin i
=\frac{2d}{\cos\gamma}\left(n_2-n_1\sin i\sin\gamma\right).
\]
结合式\eqref{eq:q1-snell}可化为
\begin{equation}
\boxed{\Delta=2n_2d\cos\gamma.}
\label{eq:q1-opd}
\end{equation}
若真空波数为$\sigma$，则真空波长为$\lambda=10^4/\sigma$。由此得到外延层内一次往返的传播相位
\begin{equation}
\Phi=\frac{2\pi\Delta}{\lambda}
=4\pi\times10^{-4}\sigma d n_2\cos\gamma,
\label{eq:q1-phase}
\end{equation}
其中$d$以$\mu\mathrm m$为单位，$10^{-4}$用于完成$\mu\mathrm m$到$\mathrm{cm}$的换算。式\eqref{eq:q1-phase}只描述外延层往返传播产生的光程相位；界面反射相变由后续带符号Fresnel振幅系数直接体现，因此总反射场中无需再附加固定$\pi$相位。

\subsubsection{Sellmeier--Drude折射率模型}
4H-SiC的折射率具有明显波长色散。Zhou等在4H-SiC红外波导色散研究中采用多项Sellmeier关系描述材料折射率\cite{zhou2023sic}；Tiwald等通过傅里叶变换红外椭偏实验测量掺氮4H/6H-SiC的介电响应，并以经典Drude项描述自由电子等离子体对红外光学响应的贡献\cite{tiwald1999sic}。据此，本文将晶格色散与自由载流子响应统一写在介电函数层面。

\textbf{（1）Sellmeier晶格色散。} SiC外延层的晶格介电函数采用三项Sellmeier形式
\begin{equation}
\varepsilon_{\mathrm L}(\lambda)
=A+\sum_{j=1}^{3}\frac{B_j\lambda^2}{\lambda^2-C_j}.
\label{eq:q1-sellmeier}
\end{equation}
其中$A,B_j$为无量纲参数，$C_j$与$\lambda^2$具有相同量纲，即单位为$\mu\mathrm m^2$。

\textbf{（2）Drude载流子修正。} 自由载流子的红外介电响应采用Drude模型\cite{drude,tiwald1999sic}
\begin{equation}
\widetilde\varepsilon_2(\omega,\lambda)
=\varepsilon_{\mathrm L}(\lambda)
-\frac{\omega_p^2}{\omega(\omega+\mathrm i\Gamma_e)},
\qquad
\omega_p^2=\frac{Ne^2}{\varepsilon_0m^*},
\label{eq:q1-drude-complex}
\end{equation}
其中$\Gamma_e$为自由载流子碰撞频率，$N$为载流子浓度，$e$为元电荷，$\varepsilon_0$为真空介电常数，$m^*$为载流子有效质量。对于SiC透明或弱吸收区域，有$\Gamma_e/\omega\ll1$，耗散虚部相对较小。

取式\eqref{eq:q1-drude-complex}的弱吸收实部近似，并利用$\omega=2\pi c/\lambda$，得到
\begin{equation}
n_2^2(\lambda,N)=\varepsilon_{\mathrm L}(\lambda)
-\frac{Ne^2\lambda^2}{4\pi^2c^2\varepsilon_0m^*}.
\label{eq:q1-drude}
\end{equation}
式\eqref{eq:q1-drude}中$N$与$m^*$以$N/m^*$的组合进入折射率，仅由单一反射光谱难以同时独立识别。Volm等通过回旋共振实验与能带计算研究了4H-SiC电子有效质量张量，理论计算给出的一个主分量约为$0.28m_0$\cite{volm1996sic}。本文将$m^*=0.28m_0$作为等效有效质量固定，仅反演载流子浓度$N$；该处理用于消除$N/m^*$的参数冗余，不将$m^*$本身作为材料常数测量结果。上述实折射率近似仅用于透明或弱吸收波段，当吸收不可忽略时应保留式\eqref{eq:q1-drude-complex}的完整复介电响应。

将式\eqref{eq:q1-drude}得到的实折射率代入式\eqref{eq:q1-snell}，即可确定折射角$\gamma$。

\subsubsection{Fresnel双光束反射率模型}
对偏振态$u\in\{s,p\}$，按照光从介质$i$传播到介质$j$的方向定义带符号Fresnel振幅反射系数$r_{ij}^{(u)}$和透射系数$t_{ij}^{(u)}$\cite{bornwolf,yeh}。其中，s偏振满足
\begin{equation}
r_{ij}^{(s)}=
\frac{n_i\cos\theta_i-n_j\cos\theta_j}
{n_i\cos\theta_i+n_j\cos\theta_j},\qquad
t_{ij}^{(s)}=
\frac{2n_i\cos\theta_i}
{n_i\cos\theta_i+n_j\cos\theta_j},
\label{eq:q1-fresnel-s}
\end{equation}
p偏振满足
\begin{equation}
r_{ij}^{(p)}=
\frac{n_j\cos\theta_i-n_i\cos\theta_j}
{n_j\cos\theta_i+n_i\cos\theta_j},\qquad
t_{ij}^{(p)}=
\frac{2n_i\cos\theta_i}
{n_j\cos\theta_i+n_i\cos\theta_j}.
\label{eq:q1-fresnel-p}
\end{equation}

表面反射光与第一束内部反射光叠加后，偏振态$u$的总复振幅反射系数为
\begin{equation}
\rho_u=r_{12}^{(u)}
+t_{12}^{(u)}r_{23}^{(u)}t_{21}^{(u)}\exp(-\mathrm i\Phi).
\label{eq:q1-double-field}
\end{equation}
在上述振幅约定下，$r_{ij}^{(u)}$的正负号已经包含界面反射所对应的相位关系。因此式\eqref{eq:q1-double-field}直接叠加界面振幅与传播相位，无需另行引入“半波损失”形式的固定$\pi$相位。对于非偏振入射光，理论反射率取s、p两种偏振反射率的算术平均：
\begin{equation}
R_{\mathrm{th}}(\sigma,i;d,N,n_3,\boldsymbol\beta)
=\frac{|\rho_s|^2+|\rho_p|^2}{2},
\label{eq:q1-reflectance}
\end{equation}
其中$\boldsymbol\beta=(A,B_1,B_2,B_3,C_1,C_2,C_3)$为Sellmeier参数向量。式\eqref{eq:q1-reflectance}由此给出厚度、色散参数和界面光学性质与可观测反射率之间的定量关系。

\subsection{模型汇总}
综上，问题一的Sellmeier--Drude--Fresnel双光束模型可整理为
\begin{equation}
\left\{
\begin{aligned}
&n_1=1.0003,\qquad \lambda=10^4/\sigma,\\
&\varepsilon_{\mathrm L}(\lambda)
=A+\sum_{j=1}^{3}\frac{B_j\lambda^2}{\lambda^2-C_j},\\
&n_2^2(\lambda,N)=\varepsilon_{\mathrm L}(\lambda)
-\frac{Ne^2\lambda^2}{4\pi^2c^2\varepsilon_0m^*},\\
&n_1\sin i=n_2\sin\gamma,\\
&\Delta=2n_2d\cos\gamma,\qquad
\Phi=4\pi\times10^{-4}\sigma d n_2\cos\gamma,\\
&\rho_u=r_{12}^{(u)}
+t_{12}^{(u)}r_{23}^{(u)}t_{21}^{(u)}e^{-\mathrm i\Phi},\quad u\in\{s,p\},\\
&R_{\mathrm{th}}=\dfrac{|\rho_s|^2+|\rho_p|^2}{2}.
\end{aligned}
\right.
\label{eq:q1-summary}
\end{equation}
给定$\sigma,i,d,N,n_3$及Sellmeier参数$\boldsymbol\beta$，式\eqref{eq:q1-summary}即可逐点计算理论反射率$R_{\mathrm{th}}$。该响应函数同时保留厚度引起的传播相位、材料色散以及界面振幅效应，问题二将在此基础上由实测光谱反演未知参数。

\subsection{模型求解}
基于式\eqref{eq:q1-summary}，构造逐波数计算理论反射率的正演算法。对每一波数，依次计算Sellmeier--Drude折射率、Snell折射角、Fresnel振幅系数与传播相位，最后对$s,p$偏振反射率取算术平均。标准伪代码如下。

\begin{table}[H]
\centering
\begin{minipage}{\textwidth}
\setlength{\parindent}{0pt}
\normalsize
\hrule
\vspace{0.45ex}
\textbf{Algorithm 1}\quad 问题一双光束理论反射率正演算法
\vspace{0.45ex}
\hrule
\vspace{0.65ex}
\textbf{Input:} 波数序列$\{\sigma_j\}_{j=1}^{M}$，入射角$i$，厚度$d$，载流子浓度$N$，衬底折射率$n_3$，Sellmeier参数$\boldsymbol\beta$；\\[-0.2ex]
\textbf{Initialize:} $n_1\leftarrow1.0003$，$m^*\leftarrow0.28m_0$，建立空反射率序列$\mathcal R$；
\vspace{0.45ex}

\begin{tabularx}{\linewidth}{@{}r@{\hspace{0.8em}}X@{}}
1: & \textbf{for} $j=1,2,\ldots,M$ \textbf{do} \\
2: & \quad 计算波长$\lambda_j\leftarrow10^4/\sigma_j$； \\
3: & \quad 由Sellmeier模型计算晶格介电函数$\varepsilon_{\mathrm L}(\lambda_j)$； \\
4: & \quad 由Drude弱吸收修正计算$n_{2,j}$；若$n_{2,j}^2\le0$，则判定该参数组无效； \\
5: & \quad 由Snell定律计算折射角$\gamma_j$； \\
6: & \quad 分别计算$s,p$偏振的$r_{12},t_{12},r_{23},t_{21}$； \\
7: & \quad 计算传播相位$\Phi_j\leftarrow4\pi\times10^{-4}\sigma_j d n_{2,j}\cos\gamma_j$； \\
8: & \quad \textbf{for} $u\in\{s,p\}$ \textbf{do} \\
9: & \qquad 计算$\rho_u\leftarrow r_{12}^{(u)}+t_{12}^{(u)}r_{23}^{(u)}t_{21}^{(u)}e^{-\mathrm i\Phi_j}$； \\
10: & \quad \textbf{end for} \\
11: & \quad 计算$R_j\leftarrow\bigl(|\rho_s|^2+|\rho_p|^2\bigr)/2$并写入$\mathcal R$； \\
12: & \textbf{end for} \\
13: & \textbf{return} $\mathcal R=\{R_j\}_{j=1}^{M}$；
\end{tabularx}

\vspace{0.45ex}
\textbf{Output:} 与输入波数逐点对应的非偏振理论反射率序列$\mathcal R$。
\vspace{0.55ex}
\hrule
\end{minipage}
\end{table}


计算过程中要求$n_2^2>0$且Snell关系存在实数解，以保证当前弱吸收模型处于其物理适用范围。由此定义正演算子
\begin{equation}
\mathcal F:\ (\sigma,i;d,N,n_3,\boldsymbol\beta)\longmapsto R_{\mathrm{th}}.
\label{eq:q1-forward-operator}
\end{equation}
算子$\mathcal F$在给定一组厚度和光学参数后生成整段理论反射光谱；问题二据此以实测反射率为目标，通过最小化理论光谱与观测光谱之间的差异识别$d$及相关参数。

\subsection{模型性质与厚度信息分析}

\subsubsection{厚度对干涉条纹的控制作用}
由式\eqref{eq:q1-phase}可直接得到
\begin{equation}
\frac{\partial\Phi}{\partial d}
=4\pi\times10^{-4}\sigma n_2\cos\gamma,
\label{eq:q1-phase-thickness-sensitivity}
\end{equation}
因此，在材料光学参数固定时，厚度$d$与传播相位严格线性。厚度增大将使相位随波数变化得更快，从而在相同波数区间内产生更密集的干涉条纹，条纹周期因而直接携带厚度信息。

进一步考虑$n_2$及$\gamma$随波数缓慢变化的情况。相邻同类干涉极值近似满足$\Delta\Phi\approx2\pi$，于是局部条纹间距可写为
\begin{equation}
\Delta\sigma(\sigma)
\approx
\frac{1}
{2\times10^{-4}d\left|\dfrac{\mathrm d}{\mathrm d\sigma}
\left[\sigma n_2(\sigma)\cos\gamma(\sigma)\right]\right|}.
\label{eq:q1-fringe-spacing-general}
\end{equation}
若在较窄波段内色散变化可以忽略，则式\eqref{eq:q1-fringe-spacing-general}退化为
\begin{equation}
\Delta\sigma
\approx\frac{1}{2\times10^{-4}d n_2\cos\gamma},
\qquad
 d\approx\frac{1}{2\times10^{-4}n_2\cos\gamma\,\Delta\sigma}.
\label{eq:q1-fringe-spacing-simple}
\end{equation}
式\eqref{eq:q1-fringe-spacing-simple}对应常折射率近似下的条纹间距测厚关系；当$n_2(\sigma)$的变化不可忽略时，式\eqref{eq:q1-fringe-spacing-general}中的色散导数将直接修正条纹间距，因此实际反演采用连续光谱和完整色散关系，而不以单一常数折射率代替$n_2(\sigma)$。

\subsubsection{色散、载流子与界面参数的作用}
厚度$d$主要通过传播相位$\Phi$控制条纹位置与密度；Sellmeier参数和载流子浓度$N$通过$n_2(\sigma)$同时影响传播相位和Fresnel系数，因而与厚度存在参数耦合。衬底折射率$n_3$不进入外延层传播相位，但会改变外延层--衬底界面的$r_{23}^{(u)}$，主要影响干涉振幅与包络；入射角$i$则同时改变折射角$\gamma$以及$s,p$偏振的界面响应。

这一作用路径与图\ref{fig:q1-flow}中的“材料光学模型--光束正演模型”关系相对应：$d$直接进入相位项，$N$和Sellmeier参数先改变$n_2$后同时传递至相位与界面系数，而$n_3$主要通过$r_{23}^{(u)}$作用于振幅。由于这些参数对整段光谱的作用形式不同，连续波段拟合能够同时利用条纹相位、局部周期和振幅包络的信息；不同入射角提供独立观测，可进一步比较厚度结果的一致性。因此后续采用分角度全谱非线性最小二乘并配合多初值搜索。

\subsection{模型检验}
问题一尚未使用实测数据反演未知参数，因此检验对象是解析关系及其数值实现，而不是拟合优度。首先利用可直接从模型推出的极限关系检查折射、相位和界面系数是否保持一致，结果见表\ref{tab:q1-selfcheck}。

\begin{table}[!htbp]
\centering
\caption{问题一模型自洽性检验}
\label{tab:q1-selfcheck}
\renewcommand{\arraystretch}{1.18}
\begin{tabularx}{0.94\textwidth}{>{\centering\arraybackslash}p{3.0cm} X >{\centering\arraybackslash}p{2.3cm}}
\toprule
\rowcolor{CumcmLightBlue}
\textbf{检验项} & \textbf{理论关系} & \textbf{结果} \\
\midrule
零载流子极限 & $N\to0$时，式\eqref{eq:q1-drude}退化为$n_2^2=\varepsilon_{\mathrm L}$ & 满足 \\
正入射极限 & $i=\gamma=0$时，式\eqref{eq:q1-fresnel-s}--\eqref{eq:q1-fresnel-p}退化为经典正入射Fresnel系数 & 满足 \\
厚度线性 & 其他参数固定时，式\eqref{eq:q1-phase-thickness-sensitivity}表明$\Phi\propto d$ & 满足 \\
相位处理 & 界面反射相变由带符号$r_{ij}^{(u)}$承担，不重复叠加固定$\pi$相位 & 一致 \\
\bottomrule
\end{tabularx}
\end{table}

\subsubsection{数值实现一致性检验}
在解析关系之外，分别计算正入射Fresnel极限、理论反射率取值范围、波数--波长单位缩放以及界面相位处理，并将程序输出与相应解析表达式逐项比较。其中传播相位的公式值与程序值最大差异为$1.25\times10^{-14}$，处于双精度浮点舍入量级；各测试点的反射率均保持有限且非负。由此可确认代码实现与式\eqref{eq:q1-snell}--\eqref{eq:q1-reflectance}采用了相同的单位体系和相位约定。

需要说明的是，$d\to0$时完整薄膜反射应退化为空气--衬底单界面反射，而当前模型按题意只截取表面反射和第一束内部反射，并非完整Airy级数，因此该极限不作为双光束截断模型的必要验证条件。

上述解析与数值检验确认了折射关系、单位换算、界面相位和程序实现的一致性。由式\eqref{eq:q1-reflectance}得到的理论反射率函数可据此用于问题二的SiC光谱参数反演。

\FloatBarrier